\section{Conclusion}
\label{sec:conclusion}

We presented the first controlled study of representation drift under self-critique in language models. By comparing residual stream activations across base, critique, and paraphrase conditions in \pythia on 150 \gsm problems, we demonstrated that self-critique induces significantly larger (41\% greater $L_2$ displacement), more layer-specific (peak at layers 12--16), and more directionally consistent (0.596 vs.\ 0.479 cosine consistency) representational drift than length-matched paraphrasing. CKA analysis confirmed that critique creates genuinely different population-level structure. These findings establish that self-critique produces measurable geometric restructuring of the residual stream, not merely surface-level variation.

Our work also identifies an important methodological caveat: linear probe accuracy on critique activations is inflated by explicit correctness information in critique text. Future studies of self-critique representations should carefully control for information content to avoid confounding evaluative structure with evaluative content.

Several directions for future work emerge. First, causal validation via activation steering along the identified critique direction would test whether this subspace functionally mediates self-evaluative behavior. Second, studying self-generated (rather than externally provided) critique would better capture natural self-reflection dynamics. Third, extending the analysis to instruction-tuned and reasoning-specialized models would establish whether critique processing is amplified by training. Finally, multi-round self-critique experiments could reveal whether representational changes compound or converge across iterations.
